\documentclass[a5paper,10pt]{article}

% https://github.com/InDevRus/filippov-solutions
% Нумерация страниц
\pagenumbering{gobble}

% Настройка полей
\usepackage{geometry}
\geometry{left=1cm}
\geometry{right=1cm}
\geometry{top=1cm}
\geometry{bottom=1.5cm}

% Вставка таблицы
\usepackage{array}
\usepackage{tabularx}

% Инструменты выравнивания формул
\usepackage{amsmath}
\usepackage{mathtools}

% Диакритика в математике
\usepackage{amsfonts}

% Вычисления размеров на ходу
\usepackage{calc}

% Удобные записи производных
\usepackage[italic=true]{derivative}

% Настройки сносок
\usepackage{enotez}

% Длинные знаки векторов
\usepackage{anyfontsize}
\usepackage[b]{esvect}

% Вставка картинок
\usepackage{graphicx}

% Вставка красной строки
\usepackage{indentfirst}
\setlength{\parindent}{1em}

% Условия в командах
\usepackage{xifthen}

% Луа-скрипты
\usepackage{luacode}

% Настройка команд отдельно в математическом и текстовом режимах
\usepackage{mathcommand}

% Для повторения LaTeX кода
\usepackage{multido}

% Конфигурация отступов формул
\delimitershortfall-1sp
\usepackage{mleftright}
\mleftright

% Растягивание объектов
\usepackage{relsize}

% Обтекание текстом
\usepackage{wrapfig}
\usepackage{subcaption}
\usepackage{floatflt}

% Поддержка ввода кириллицы
\usepackage[english, russian]{babel}

% Настройка корневой позиции для компилятора
\usepackage{import}
\providecommand*{\ProjectRootPath}{..}

\import{\ProjectRootPath/preambles/macros/}{theorems}

\import{\ProjectRootPath/preambles/fonts}{font_setup}

% Знак градуса
\usepackage{siunitx}

\import{\ProjectRootPath/preambles/macros/}{operators}
\import{\ProjectRootPath/preambles/macros/}{redefinitions}

\import{\ProjectRootPath/preambles/fonts}{letters_setup}
\import{\ProjectRootPath/preambles/fonts/adjustments}{custom_superscripts}

\import{\ProjectRootPath/preambles/custom}{formatting}
\import{\ProjectRootPath/preambles/custom}{frames}
\import{\ProjectRootPath/preambles/custom}{list_setup}

% TODO: перевести команды на современный стиль объявления
\input{../preambles/plotting}

\begin{document}

$\zed'\br{t} = \tpow{10}{\hspace{0.5mm}x + \zed}$.
$\tpow{10}{-\zed} \cdot \zed' = \tpow{10}{\hspace{0.5mm}x}$.
\vspace{1mm}
Взяв интеграл левой и правой частей по $dx$, получаем $-\dfrac {\tpow{10}{-\zed}}{\ln 10} = \dfrac{\tpow{10}{\hspace{0.5mm}x}} {\ln 10} + C$.

Окончательно, общее решение равняется
$\zed\br{x} = -\lg\br{C - \tpow{10}{\hspace{0.5mm}x}}$, где
$C > 0$. Область определения: $x \in \br{-\infty; \lg\br{C}}$.

\begin{figure}[!h]
    \centering
    \begin{tikzpicture}
        \pgfplotsset{
            Graph/.style = {
                samples = 200,
            },
        }
        \begin{axis}[
            width = 12cm,
            ylabel = $\zed$,
            ymin = -4.2,
            ymax = 4.2,
            xmin = -5.2,
            xmax = 5.2, 
            xtick = {-10, ..., 10},
            ytick = {-10, ..., 10},
            minor tick num = 3,
            declare function = {
                f(\C, \x) = -log10(\C - 10^(x));
                g(\C, \x) = log10(\C - 10^(-x));
            }]
            
            \draw[thick, blue] (-5.5, 0) -- (5.5, 0);
            
            \pgfplotsinvokeforeach {0.001, 0.003, 0.01, 0.03, 0.1, 0.25, 0.5, 1, 3, 11, 32, 100, 300, 1000, 3000, 10000} {
                \addplot[Graph, domain = -5.5 : log10(#1 / 2)]
                {f(#1, x)};
                \addplot[Graph, domain = -log10(#1 / 2) : 5.5]
                ({g(#1, x)}, {x});
            }
        \end{axis}
    \end{tikzpicture}
\end{figure}

\end{document}
